\begin{table*}[t]
\centering
\small
\resizebox{\textwidth}{!}{%
\begin{tabular}{llcccccccr}
\toprule
\textbf{Selection Strategy} & \textbf{Scoring Signal} & \textbf{Passage Coverage} & \textbf{Budget $B$} & \textbf{Draft Tokens (Comp)} & \textbf{$\alpha_{\text{eval}}$ (\%)} & \textbf{Tok/Rnd} & \textbf{E2E ms/tok} & \textbf{E2E tok/s} & \textbf{Speedup vs. Full} \\
\midrule
BM25 Sentences & Lexical BM25 & Yes ($m/m$) & 10,000 & 1,646.4 (-67.7\%) & 77.24\% & 3.047 & 23.45 & 42.65 & +2.8\% \\
\textbf{SPADE (Proposed)} & Target Attention Hook & \textbf{Yes ($m/m$)} & 10,000 & 1,776.8 (-65.1\%) & \textbf{78.20\%} & \textbf{3.114} & \textbf{23.10} & \textbf{43.28} & \textbf{+4.4\%} \\
\bottomrule
\end{tabular}%
}
\caption{\textbf{Ablation of Evidence Selection Strategy and Coverage Rule:} Comparing uniform random, lexical BM25, and target-attention sentence selection under an identical budget ($B=10\text{k}$ characters) on 300 held-out claims. Disabling passage coverage guarantees causes cited document omission and proposal divergence, while attention scoring with passage-level coverage achieves the highest acceptance ($78.20\%$) and maximum wall-clock speedup ($+4.4\%$ vs.\ Full).}
\label{tab:selector_ablation}
\end{table*}
